%---------------------------------------------------------------
% SLO: slovenski povzetek
% ENG: slovenian abstract
%---------------------------------------------------------------
\selectlanguage{slovene} % Preklopi na slovenski jezik
\addcontentsline{toc}{chapter}{Povzetek}
\chapter*{Povzetek}

\noindent\textbf{Naslov:} \ttitle
\bigskip

Sestava naložbenega portfelja je klasičen optimizacijski problem, ki ga je Markowitz formuliral kot iskanje ravnovesja med pričakovanim donosom in tveganjem. V praksi pri tem naletimo na dve težavi: pričakovane donose je težko oceniti, realistična kardinalna omejitev na število sredstev v portfelju pa problem spremeni v mešano celoštevilsko in nekonveksno nalogo, ki je klasične kvadratne metode ne rešujejo. V magistrskem delu predstavljamo in odprtokodno implementiramo hibridni cevovod, ki obe težavi naslavlja ločeno. Pričakovane donose in kovariančno matriko ocenjujemo s presečnim (angl.\ \emph{cross-sectional}) transformerjem tipa MASTER, ki iz cenovne zgodovine vseh sredstev hkrati napove njihovo razvrstitev po prihodnjem donosu ter faktorsko kovarianco iz iste deljene predstavitve; napovedi kot dodatno plast robustnosti sproti spletno kombiniramo z zgodovinskim povprečjem po postopku Hedge. Kardinalno omejeni problem nato rešimo s tremi metahevristikami v jeziku C++ (rojenje delcev, simulirano ohlajanje, genetski algoritem). Kakovost optimizatorjev neodvisno ovrednotimo na knjižnici OR-Library, celoten cevovod pa s postopnim (walk-forward) testiranjem v šestih tržnih režimih ter z dvema neposrednima primerjavama na protokolu objavljenih študij. Transformerski cevovod prekaša klasičnega v petih od šestih režimov, je edini vir napovedi z zunaj-vzorčno pozitivno napovedno močjo in edini (skupaj z ansamblom) prestane na izbiro popravljeni Sharpov količnik; prednost pred golim zgodovinskim povprečjem ostaja ekonomsko smiselna, a statistično mejna.

\subsection*{Ključne besede}
\textit{\tkeywords}
\clearemptydoublepage